\section{Case III: The Gauge That Did Not Transfer}
\label{sec:gauge}

A recurrent composer is trained on the word problems of five finite
groups and scored for compositional depth generalization: trained at
depth $\le8$, tested at depth 64. Two
pre-registered lenses score recovery of the group representation from
the state: the \emph{primary} de-gauges with a scale-only correction,
valid when the learned basis matches the target basis up to scale
($\hat{Q}\approx I$), measured to hold on every checkpoint of the
\emph{previous} architecture; the \emph{crosscheck}
fits a full orthogonal intertwiner by Procrustes on a fit split and
scores on a disjoint eval split, assuming nothing about the gauge.

\looseness=-1 \textbf{The signature.} The 62-cell sweep completed cleanly, and the
endpoint machinery, run as built on the primary lens, returned
\textsc{falsify}: no measurable separation from the baseline at the
decisive groups.
% <!-- evidence: X1 -->
The mandatory instrument-health pass (Rule 1) then surfaced a
contradiction (Figure~\ref{fig:lens}): wherever training had
\emph{not} converged, both lenses read near zero; wherever it
converged, they contradicted each other at full magnitude (one
converged cell: final loss $0.0001$, primary recovery $0.050$,
crosscheck $1.000$). All 13 cells with a lens disagreement of at least
$0.5$ are converged cells; on all 30 baseline cells the lenses agree
exactly. The harvest stopped for adjudication with no model verdict
recorded.
% <!-- evidence: X2 -->

\looseness=-1 \textbf{The tiebreak.} A leaky crosscheck reading high on junk would
leave the falsification standing, so either lens could be the broken
one. The adjudication, performed against raw artifacts and code
(Rule 4), rested on four grounds:
(1)~the crosscheck is leakage-guarded by construction (fit on the fit
split, scored on the disjoint eval split); (2)~it
demonstrably discriminates, reading near zero on every non-converged
cell; (3)~the primary's scale-only correction was already flagged as
basis-brittle across seeds (per-seed mean cosine $0.03$ to $0.69$
where the crosscheck read $0.86$ to $0.95$), and its
$\hat{Q}\approx I$ evidence came from a different architecture; a
perfect model in a rotated basis reads zero under an identity gauge by
construction;
% <!-- evidence: X3 -->
(4)~the synthetic acceptance test already grades the crosscheck when
the injected gauge is non-identity (Rule 2).
Verdict: the primary lens is broken on converged cells; the mechanical
falsification is void \emph{as a model verdict}.

\looseness=-1 \textbf{Teeth.} The tiebreak carried its own pre-registered falsifier
(Rule 3): re-scoring with shuffled targets must read near zero, or the
tiebreak is wrong and the contradiction escalates to an instrument
rebuild. On three converged checkpoints the shuffled crosscheck read
$0.00$, $0.00$, and $0.05$, far below the $0.5$ void threshold, while
the unshuffled reads reproduced the committed values bit-identically
(Table~\ref{tab:teeth}).
% <!-- evidence: X4 -->

\textbf{The re-metric.} The full grid was re-read under the crosscheck
lens through the unmodified verdict machinery, which first reproduced
the committed primary verdict and all 62 recomputed ceilings to within
$10^{-9}$.
% <!-- evidence: X5 -->
The corrected endpoint is \emph{still} \textsc{falsify}, for a more
informative reason (\S\ref{sec:notflipped}, Table~\ref{tab:remetric}):
not ``no signal anywhere'' but ``two specific, lens-independent
failures at the two decisive groups.'' Fixing the instrument changed
the reason, not the verdict.

\begin{figure}[t]
\centering
\includegraphics[width=0.48\linewidth]{fig1_lens_contradiction}
\caption{Primary versus crosscheck recovery (fraction of items
recovered at cosine $\geq0.9$, depth 64), all 62 sweep cells.
Non-converged cells (open orange) sit near the origin under both
lenses; every cell with a lens disagreement of at least $0.5$ is a
converged cell (filled blue). A lens defect that activates exactly
where the model succeeds is the inverse of noise.}
\label{fig:lens}
\end{figure}
